\section{Checks supporting execution and experience}
\label{sec:checks}
Tool use and learning need records of what happened. The execution controller binds
checks to the submitted code, inputs and numerical claims, recording their digests
and results. Changed artifacts invalidate earlier receipts. A required-check policy
reports \texttt{PASS} when all required checks ran and passed, \texttt{FAIL} for a
detected defect, and otherwise \texttt{INCOMPLETE} for missing or rejected evidence.
This mechanism establishes the scope of execution evidence; it does not choose the
best method for the task.

For learning, availability and accounting checks constrain which delayed outcomes
can be consumed. Pending feedback waits and invalid feedback is rejected; valid
negative outcomes remain eligible. The purpose is to support interpretable experience
updates, not to retain only profitable episodes. Passing these checks does not prove
that supplied records are authentic or that the original choice was optimal.

The supporting audit study contrasts guidance, optional checking and required final
checking. Its accepted model submissions were not independently gradable, so it does
not establish an improvement in decision quality. Appendix~\ref{app:audit} retains
the protocol, historical pilot, instrument failures and regression evidence. Those
records define the current checking scope separately from the capability study.
